\section{Question 2: Model Compression Implementation}

\subsection{(a) Configurable compression method}

\textbf{Choice of method.} I use quantization-aware training (QAT) with a
configurable mixed-precision policy. QAT is preferred over ordinary
post-training quantization (PTQ) for the target low precisions because PTQ
only calibrates a completed FP32 model, whereas QAT lets the weights and
quantization ranges adapt to rounding and clipping error during fine-tuning.
This is especially relevant to MobileNetV2: its depthwise layers, narrow
linear bottlenecks, and residual additions make it sensitive to poorly chosen
scales. Jacob \emph{et al.} report that training co-designed with quantization
improves the accuracy--latency trade-off for MobileNets
\cite{jacob2018quantization}; LSQ and PACT further show that learned step
sizes and learned activation clipping are effective at low precision
\cite{esser2020lsq,choi2018pact}. PTQ is nevertheless retained in this
repository as a diagnostic for checking quantizer coverage and identifying
range-sensitive layers before QAT.

\paragraph{Quantizer used during QAT.}
For a value $v$, positive scale $s$, and integer bounds $Q_n,Q_p$, the
forward pass uses
\[
 q = \operatorname{clip}\!\left(\operatorname{round}(v/s),Q_n,Q_p\right),
 \qquad \widehat v=sq.
\]
Here $\widehat v$ is the dequantized FP32 value consumed by the next PyTorch
operation. Rounding uses a straight-through estimator (STE): it is discrete in
the forward pass but has an identity surrogate gradient in the backward pass.
Thus, the model learns under the error that integer rounding and clipping would
introduce, while the master weights and optimizer state remain FP32.

\begin{itemize}
    \item \textbf{Weights:} every convolution and linear weight tensor uses
    symmetric signed quantization with
    $Q_n=-(2^{b-1}-1)$ and $Q_p=2^{b-1}-1$. One learned FP32 scale is stored
    for each output channel. The scale is initialized using the LSQ rule and
    receives LSQ-style gradient scaling.
    \item \textbf{ReLU6 outputs:} these are non-negative and bounded, so they
    use unsigned levels $[0,2^b-1]$, one learned scale per tensor, and initial
    scale $6/(2^b-1)$. This is a PACT-like learned clipping-range treatment.
    \item \textbf{Signed activation boundaries:} normalized input,
    linear-projection outputs, residual operands/results, and logits can be
    negative. They therefore use a symmetric signed, learned per-tensor scale.
    \item \textbf{Residual addition:} the skip and main-branch operands are
    fake-quantized with the \emph{same} signed scale before addition; the sum
    is quantized again with that boundary scale. This prevents invalid integer
    addition of tensors represented by unrelated scales.
    \item \textbf{Precision schedule:} a QAT candidate progresses from W8/A8
    to W6/A6 and then to its requested target. The final target policy is
    validation-selected only after it has been reached.
\end{itemize}

\paragraph{Training and export procedure.}

\begin{itemize}
    \item Each QAT run begins from the same immutable FP32 baseline checkpoint.
    The implementation fine-tunes for 12 epochs using SGD with Nesterov
    momentum, cosine decay, and a separate zero-weight-decay parameter group
    for learned quantizer scales.
    \item BatchNorm remains FP32 during the adaptation phase, then its running
    statistics and affine parameters are frozen from epoch 10.
    \item The best target-precision checkpoint is selected on the fixed
    5,000-image validation split. The test set is not used to select a bit
    policy.
    \item At export, each adjacent Conv--BatchNorm pair is folded. The
    exporter packs integer weight codes and writes the scales, int32 biases,
    requantization fields, descriptors, alignment padding, and header into a
    custom \texttt{QPK1} artifact.
\end{itemize}

\subsection{(b) Applying QAT to MobileNetV2}

The wrappers in \texttt{src/quantization.py} and
\texttt{src/compression.py} replace all eligible \texttt{Conv2d},
\texttt{Linear}, and \texttt{ReLU6} modules. They also add explicit
quantizers at graph boundaries that are not represented by a standalone
torchvision layer. The exact policy is summarized in
Table~\ref{tab:q2-coverage}.

\begin{table}[H]
    \centering
    \small
    \caption{QAT coverage for MobileNetV2. The selected compact policy is W4/A6 with W6 depthwise weights and W8 stem/classifier weights.}
    \label{tab:q2-coverage}
    \begin{tabular}{p{0.27\linewidth}p{0.39\linewidth}p{0.22\linewidth}}
        \toprule
        Layer or boundary & Fake-quantization operation & Selected precision \\
        \midrule
        Normalized image input and final logits & Learned per-tensor signed activation quantizer & A6 \\
        Stem convolution and classifier & Per-output-channel signed weight quantizer; convolution/linear is evaluated on dequantized FP32 values during QAT & W8 weights \\
        Expansion and projection $1\times1$ convolutions & Per-output-channel signed weight quantizer; a non-residual projection output receives a signed activation quantizer & W4 weights; A6 boundary \\
        Depthwise $3\times3$ convolutions & Per-output-channel signed weight quantizer; no depthwise layer is omitted & W6 weights \\
        ReLU6 & ReLU6 followed by learned per-tensor unsigned fake quantization & A6 \\
        Residual addition & Fake-quantize skip and branch with one shared signed scale, add, then fake-quantize the result & A6 \\
        BatchNorm & Retained in FP32 for QAT adaptation; folded into the preceding convolution for export & Folded; no separate stored BN \\
        Pooling, flatten, and dropout & No weights are quantized; dropout is disabled at inference & No stored payload \\
        \bottomrule
    \end{tabular}
\end{table}

The policy is configurable from the command line. The default bit width is
applied to all quantized convolution and linear layers; optional controls can
raise only depthwise weights or the first and last weighted layers. Exact named
overrides are checked against the model graph, so a misspelled layer cannot be
silently skipped. A coverage audit records each weighted module's name, shape,
bit width, scale shape, signedness, and exception reason.

\subsection{(c) Storage overheads and deployment limitations}

\paragraph{Storage accounting.} The reported size is the deployment artifact,
not a PyTorch checkpoint. The checkpoint is intentionally larger because it
contains FP32 master weights, optimizer state, and QAT state needed to resume
training. The \texttt{QPK1} exporter instead charges all inference-time
components:

\begin{itemize}
    \item packed signed weight codes: $\lceil Nb/8\rceil$ bytes for $N$
    weights at $b$ bits;
    \item one FP32 scale per output channel for each convolution or linear
    layer, plus one FP32 scale per activation boundary;
    \item int32 folded biases and one int32 multiplier/int32 shift pair per
    output channel for requantization;
    \item compact JSON descriptors, a 12-byte header, and 16-byte alignment
    padding.
\end{itemize}

For the selected W4/A6 mixed-precision artifact, ordinary weights are W4,
depthwise weights are W6, and stem/classifier weights are W8. The byte-verified
\texttt{deployable\_model.qpk} is 1,416,800 bytes (1.351 MiB), including
292,632 bytes of metadata, versus 8,878,504 bytes for its folded FP32 weighted
layer reference. This gives a 6.267$\times$ model-size reduction. The
repository verifies that this analytical total equals the physical artifact
length.

Activation storage is measured for one $1\times3\times32\times32$ inference
by tracing quantizer-boundary liveness. Residual operands remain live until
their corresponding add. This produces 393,216 FP32 peak-live bytes versus
73,732 quantized peak-live bytes; boundary traffic is 2,334,248 FP32 bytes
versus 437,968 quantized bytes.

\paragraph{PyTorch limitation.} This implementation is \emph{fake QAT}, not
a benchmark of true integer inference. Although the forward pass simulates
the intended low-bit codes through fake quantization, PyTorch's
\texttt{conv2d} and \texttt{linear} still execute FP32 kernels on the
dequantized tensors. The custom \texttt{QPK1} file contains packed codes and
integer-oriented metadata, but this repository does not yet include an integer
convolution runtime that consumes it. Consequently, the reported accuracy,
storage, and activation-memory accounting are valid for the quantized model
representation; no latency, throughput, or energy speed-up is claimed.

\begin{thebibliography}{9}
\bibitem{jacob2018quantization}
B. Jacob, S. Kligys, B. Chen, M. Zhu, M. Tang, A. Howard, H. Adam, and
D. Kalenichenko, Quantization and Training of Neural Networks for Efficient
Integer-Arithmetic-Only Inference, \emph{CVPR}, 2018.
\href{https://arxiv.org/abs/1712.05877}{arXiv:1712.05877}.

\bibitem{esser2020lsq}
S. K. Esser, J. L. McKinstry, D. Bablani, R. Appuswamy, and D. S. Modha,
Learned Step Size Quantization, \emph{ICLR}, 2020.
\href{https://openreview.net/forum?id=rkgO66VKDS}{OpenReview: rkgO66VKDS}.

\bibitem{choi2018pact}
J. Choi, Z. Wang, S. Venkataramani, P. I.-J. Chuang, V. Srinivasan, and
K. Gopalakrishnan, PACT: Parameterized Clipping Activation for Quantized
Neural Networks, 2018.
\href{https://arxiv.org/abs/1805.06085}{arXiv:1805.06085}.
\end{thebibliography}
